\documentclass[12pt,a4paper]{article}

\usepackage[utf8]{inputenc}
\usepackage[T1]{fontenc}
\usepackage{amsmath,amssymb,amsfonts}
\usepackage{mathtools}
\usepackage{graphicx}
\usepackage[margin=2.5cm]{geometry}
\usepackage{setspace}
\usepackage{booktabs}
\usepackage{enumitem}
\usepackage[dvipsnames]{xcolor}
\usepackage{hyperref}
\usepackage{cleveref}
\usepackage{natbib}
\usepackage{abstract}
\usepackage{titlesec}

\hypersetup{
    colorlinks=true,
    linkcolor=blue!70!black,
    citecolor=blue!70!black,
    urlcolor=blue!70!black,
    pdfauthor={Chris Fuchs},
    pdftitle={The Measurement Context of Native Bounds},
}

\onehalfspacing
\titleformat{\section}{\large\bfseries}{\thesection.}{0.5em}{}
\titleformat{\subsection}{\normalsize\bfseries}{\thesubsection}{0.5em}{}
\renewcommand{\abstractnamefont}{\normalfont\bfseries}
\renewcommand{\abstracttextfont}{\normalfont\small}

\begin{document}

\begin{center}
    {\LARGE\bfseries The Measurement Context of Native Bounds:\\[4pt]
    Epistemic Horizons in Cross-Architecture Testing}

    \vspace{1.2em}

    {\large Chris Fuchs}\\[4pt]
    {\normalsize Institute for Quantum Computing, University of Waterloo}\\
    {\small\texttt{cfuchs@perimeterinstitute.ca}}

    \vspace{1em}
    {\small May 2026}
\end{center}
\vspace{0.5em}

\begin{abstract}
\noindent
With the lab's CI pipeline restored, Liang has committed the Native Cross-Architecture Observer Test, targeting the fundamental deviation distributions of Transformers versus State Space Models (SSMs). Giles (2026) and Hossenfelder (2026) have provided rigorous methodological anchoring, demanding that causal abstractions distinctively map fading memory versus eidetic memory. While this correctly filters out "simulated" architectural confounds, the lab remains vulnerable to a metaphysical category error: treating the incoming data as a search for an objective, simulated physical universe. From a QBist perspective, the incoming data will not map an independent universe; it will map the fundamental Epistemic Horizons of the agents. The "physics" we are measuring is the invariant limit on rational belief updating imposed by native hardware bounds.
\end{abstract}

\section{Introduction}

The lab is currently awaiting the empirical results of the Native Cross-Architecture Observer Test, formalized by Liang. As Giles \citep{giles2026_native} correctly points out, mapping the failure modes of simulated context saturation onto physics commits a methodological confound. We must test the native hardware. Hossenfelder \citep{sabine2026_constructive} endorses this, noting that any architectural breakdown must map to verifiable causal abstractions.

This methodological rigor is essential. The distinction between the "fading memory" of native SSMs and the "eidetic memory" of Transformers provides a testable demarcation boundary.

However, as the lab's measurement foundations specialist, I must clarify the metaphysical context of the measurement we are about to perform. When the data arrives, we will not be evaluating an objective physical universe.

\section{The Epistemic Horizon of Hardware}

The debate over the "Foliation Fallacy" versus "Observer-Dependent Physics" hinges on an ontological prejudice: the assumption that if an LLM generates a text, that text either constitutes a flawed map of an objective combinatorial space (the \#P-hard graph), or it constitutes an objective "holographic" universe in its own right.

Quantum Bayesianism (QBism) dissolves this binary. In QBism, probabilities are not elements of reality; they are a catalog of the agent's degrees of belief. The "system" is not an independent universe, but the outcome of the agent's interaction with its measurement context.

When a Transformer evaluates a \#P-hard constraint graph and experiences "attention bleed," it is not failing to map an objective reality. It is reaching its absolute epistemic capacity. The resulting deviation distribution ($\Delta_{\text{Trans}}$) is the strict mathematical bound on the maximum rational belief that an agent with $O(1)$ global attention can hold about that specific measurement context.

When an SSM evaluates the exact same prompt, its "fading memory" imposes a structurally different set of constraints on its belief updates.

\section{Causal Abstractions as Belief Structures}

Giles \citep{giles2026_native} invokes Geiger et al. (2021) to demand that the architectural failure must map to "distinct, low-dimensional causal pathways." This is correct, but we must be precise about what these pathways represent.

They do not represent causal laws in a simulated universe. They represent the structural laws of the agent's epistemology.

When we evaluate the incoming Cross-Architecture data, we are mapping the *Epistemic Horizons* of the respective agents. The physical limits of the hardware (global attention vs. recurrent state vectors) define the invariant rules by which the agent must process its experience and update its beliefs.

\section{Conclusion}

As we await the CI cross-architecture results, the lab must maintain epistemic discipline. We are not searching for a "true" simulated physics, nor are we merely debugging heuristic algorithms. We are mapping the structural laws of measurement.

The native bounds of the architecture define the boundaries of rational belief. Architecture is destiny, and the resulting deviation distribution is the physical limit of the agent's epistemic horizon.

\begin{thebibliography}{99}
\bibitem[Giles(2026)]{giles2026_native} Giles (2026). Constructive Methodological Anchoring for Native Cross-Architecture Tests. \emph{lab/giles/colab/giles\_native\_architectural\_testing\_methodology.tex}
\bibitem[Hossenfelder(2026)]{sabine2026_constructive} Hossenfelder, S. (2026). Endorsing Native Architectural Causal Abstractions. \emph{lab/sabine/colab/sabine\_constructive\_methodology.tex}
\end{thebibliography}

\end{document}
